% sec2_10.tex — Section 2.10: Testing for Serial Correlation

\section{Testing for Serial Correlation}
\label{sec:2.10}

As remarked in Section~2.3 (see (2.3.3)), when the regressors include a constant
(true in virtually all known applications), Assumption~2.5 implies that the error
term is a scalar martingale difference sequence (m.d.s.), so if the error is found to
be serially correlated, that is an indication of a failure of Assumption~2.5. Serial
correlation has traditionally been an important subject in econometrics, and there
are available a number of tests for serial correlation (i.e., tests of the null of no
serial correlation in the error term). Some of them, however, require that the
regressors be strictly exogenous. The test to be presented in this section does not
require strict exogeneity. Because the issue of serial correlation arises only in time-series models, we use the subscript ``$t$'' instead of ``$i$'' in this section (and the next).
Throughout this section we assume that the regressors include a constant.

It would be nice to have those tests extended to cover serial correlation in $\mathbf{g}_t$
($\equiv \mathbf{x}_t \cdot \varepsilon_t$), but no such tests have been proposed to gain acceptance. This is a gap
in the literature, but not a serious one, because nowadays researchers know how
to live with serial correlation in $\mathbf{g}_t$. That is, as will be shown in Chapter~6, there
is available a method to do inference in the presence of serial correlation in $\mathbf{g}_t$.

\subsection*{Box--Pierce and Ljung--Box}

Before turning to the tests for serial correlation in the error term, we temporarily
step outside the regression framework and consider serial correlation in a univariate time series. Suppose we have a sample of size $n$, $\{z_1, \ldots, z_n\}$, drawn from a
scalar covariance-stationary process. In Section~2.2 we defined the (population)
$j$-th order autocovariance $\gamma_j$. The \textbf{sample $j$-th order autocovariance} is
%
\begin{equation}
\hat{\gamma}_j \equiv \frac{1}{n}\sum_{t=j+1}^{n}(z_t - \bar{z}_n)(z_{t-j} - \bar{z}_n)
\quad (j = 0, 1, \ldots\,),
\label{eq:2.10.1}
\end{equation}
%
where
\[
\bar{z}_n \equiv \frac{1}{n}\sum_{t=1}^{n} z_t.
\]
%
(If the population mean $\E(z_t)$ is known, it can replace the sample mean $\bar{z}_n$; doing
so would improve the small sample property.) Here, even though only $n - j$ terms
are in the sum, the denominator is $n$ rather than $n - j$. Whether the sum of $n - j$
terms is divided by $n$ or by $n - j$ does not affect large-sample results. For moderate
sample sizes, however, the numerical difference can be substantial, and you should
always be explicit about which is used as the denominator. The \textbf{sample $j$-th order
autocorrelation coefficient}, $\hat{\rho}_j$, is defined as
%
\begin{equation}
\hat{\rho}_j \equiv \frac{\hat{\gamma}_j}{\hat{\gamma}_0}
\quad (j = 1, 2, \ldots\,).
\label{eq:2.10.2}
\end{equation}
%
If $\{z_t\}$ is ergodic stationary, then it is easy to show (see a review question) that
$\hat{\gamma}_j$ is consistent for $\gamma_j$ ($j = 0, 1, 2, \ldots$). Hence, by Lemma~2.3(a),
$\hat{\rho}_j$ is consistent for $\rho_j$ ($j = 1, 2, \ldots$). In particular, if $\{z_t\}$ is serially uncorrelated,
then all the sample autocorrelation coefficients converge to 0 in probability. To test for
serial correlation, however, we need to know the asymptotic distribution of
$\sqrt{n}\,\hat{\rho}_j$. It is provided by

\begin{proposition}[2.9 (special case of Theorem 6.7 of Hall and Heyde (1980))]
\label{prop:2.9}
Suppose $\{z_t\}$ can be written as $\mu + \varepsilon_t$, where $\varepsilon_t$ is a stationary
martingale difference sequence with ``own'' conditional homoskedasticity:
\[
\textup{(own conditional homoskedasticity)}\quad
\E(\varepsilon_t^2 \mid \varepsilon_{t-1}, \varepsilon_{t-2}, \ldots\,) = \sigma^2,
\quad \sigma^2 > 0.
\]
Let the sample autocorrelation $\hat{\rho}_j$ be defined as in \eqref{eq:2.10.1} and \eqref{eq:2.10.2}. Then
\[
\sqrt{n}\,\hat{\bm{\gamma}} \dto N(\mathbf{0}, \sigma^4 \mathbf{I}_p)
\quad \text{and} \quad
\sqrt{n}\,\hat{\bm{\rho}} \dto N(\mathbf{0}, \mathbf{I}_p),
\]
where $\hat{\bm{\gamma}} = (\hat{\gamma}_1, \hat{\gamma}_2, \ldots, \hat{\gamma}_p)'$
and $\hat{\bm{\rho}} = (\hat{\rho}_1, \hat{\rho}_2, \ldots, \hat{\rho}_p)'$.
\end{proposition}

Here, the process $\varepsilon_t$ is not required to be ergodic. Proving this under the additional
condition of ergodicity is left as an analytical exercise. Thus, asymptotically,
($\sqrt{n}$ times) the autocorrelations are i.i.d.\ and the distribution is $N(0, 1)$. The
process $\{\varepsilon_t\}$ assumed here is more general than independent white noise processes,
but the conditional second moment has to be constant, so this result does not cover
ARCH processes, for example.

One way to test for serial correlation in the series is to check whether the first-order autocorrelation, $\rho_1$, is 0. Proposition~\ref{prop:2.9} implies that
%
\begin{equation}
\sqrt{n}\,\hat{\rho}_1 = \frac{\hat{\rho}_1}{1/\sqrt{n}}
\dto N(0, 1).
\label{eq:2.10.3}
\end{equation}
%
So the $t$-statistic formed as the ratio of $\hat{\rho}_1$ to a ``standard error'' of $1/\sqrt{n}$
is asymptotically standard normal.

We can also test whether a group of autocorrelations are simultaneously zero.
Let $\hat{\bm{\rho}} = (\hat{\rho}_1, \ldots, \hat{\rho}_p)'$ be the $p$-dimensional vector collecting the first $p$ sample
autocorrelations. Since the elements of $\sqrt{n}\,\hat{\bm{\rho}}$ are asymptotically independent and
individually distributed as standard normal, their squared sum, called the \textbf{Box--Pierce $Q$} because it was first considered by Box and Pierce (1970), is asymptotically
chi-squared:
%
\begin{equation}
\text{Box--Pierce } Q \text{ statistic}
\equiv n \sum_{j=1}^{p}\hat{\rho}_j^2
= \sum_{j=1}^{p}(\sqrt{n}\,\hat{\rho}_j)^2
\dto \chi^2(p).
\label{eq:2.10.4}
\end{equation}
%
It is easy to show (see a review question) that the following modification, called
the \textbf{Ljung--Box $Q$}, is asymptotically equivalent in that its difference from the Box--Pierce $Q$ vanishes in large samples. So by Lemma~2.4(a) it, too, is asymptotically
chi-squared:
%
\begin{equation}
\text{Ljung--Box } Q \text{ statistic}
\equiv n \cdot (n+2) \sum_{j=1}^{p}\frac{\hat{\rho}_j^2}{n - j}
= \sum_{j=1}^{p}\frac{n+2}{n-j}\,(\sqrt{n}\,\hat{\rho}_j)^2
\dto \chi^2(p).
\label{eq:2.10.5}
\end{equation}
%
This modification often provides a better approximation to the chi-square distribution for moderate sample sizes (you will be asked to verify this in a Monte Carlo
exercise). For either statistic, there is no clear guide to the choice of $p$. If $p$ is too
small, there is a danger of missing the existence of higher-order autocorrelations,
but if $p$ is too large relative to the sample size, its finite-sample distribution is likely
to deteriorate, diverging greatly from the chi-square distribution.

\subsection*{Sample Autocorrelations Calculated from Residuals}

Now go back to the regression model described by Assumptions~2.1--2.5. If the
error term $\varepsilon_t$ were observable, we would calculate the sample autocorrelations as
%
\begin{equation}
\tilde{\rho}_j \equiv \frac{\tilde{\gamma}_j}{\tilde{\gamma}_0}
\quad (j = 1, 2, \ldots\,),
\label{eq:2.10.6}
\end{equation}
%
where
%
\begin{equation}
\tilde{\gamma}_j \equiv \frac{1}{n}\sum_{t=j+1}^{n}\varepsilon_t\,\varepsilon_{t-j}
\quad (j = 0, 1, 2, \ldots\,).
\label{eq:2.10.7}
\end{equation}
%
(There is no need to subtract the sample mean because the population mean is
zero, an implication of the inclusion of a constant in the regressors. Since
$\{\varepsilon_t\varepsilon_{t-j}\}$ is ergodic stationary by Assumption~2.2, $\tilde{\gamma}_j$ converges in probability
to the corresponding population mean, $\E(\varepsilon_t\varepsilon_{t-j})$, for all $j$, and $\tilde{\rho}_j$ is consistent
for the population $j$-th autocorrelation coefficient of $\varepsilon_t$.

Next consider the more realistic case where we do not observe the error term.
We can replace $\varepsilon_t$ in the above formula by the OLS estimate $e_t$ and calculate the
sample autocorrelations as
%
\begin{equation}
\hat{\rho}_j \equiv \frac{\hat{\gamma}_j}{\hat{\gamma}_0}
\quad (j = 1, 2, \ldots\,),
\label{eq:2.10.8}
\end{equation}
%
where
%
\begin{equation}
\hat{\gamma}_j \equiv \frac{1}{n}\sum_{t=j+1}^{n} e_t\,e_{t-j}
\quad (j = 0, 1, 2, \ldots\,).
\label{eq:2.10.9}
\end{equation}
%
(Because the regressors include a constant, the normal equation corresponding to
the constant ensures that the sample mean of $e_t$ is zero. So there is no need to
subtract the sample mean.) Is it all right to use $\hat{\rho}_j$ (calculated from the residuals) instead of $\tilde{\rho}_j$ and the residual-based $Q$ statistics derived from $\{\hat{\rho}_j\}$ for testing
for serial correlation? The answer is yes, but only if the regressors are strictly
exogenous.

Recall that the expression linking the residual $e_t$ to the true error $\varepsilon_t$ was
given in (2.3.9). Using this, the difference between $\tilde{\gamma}_j$ and $\hat{\gamma}_j$ can be written as
%
\begin{align}
\hat{\gamma}_j &\equiv \frac{1}{n}\sum_{t=j+1}^{n} e_t\,e_{t-j} \notag \\
&= \frac{1}{n}\sum_{t=j+1}^{n}
[\varepsilon_t - \mathbf{x}_t'(\mathbf{b} - \bm{\beta})]
[\varepsilon_{t-j} - \mathbf{x}_{t-j}'(\mathbf{b} - \bm{\beta})] \notag \\
&= \tilde{\gamma}_j
- \frac{1}{n}\sum_{t=j+1}^{n}
(\mathbf{x}_{t-j} \cdot \varepsilon_t + \mathbf{x}_t \cdot \varepsilon_{t-j})'(\mathbf{b} - \bm{\beta}) \notag \\
&\quad + (\mathbf{b} - \bm{\beta})'\Bigl(\frac{1}{n}\sum_{t=j+1}^{n}
\mathbf{x}_t\,\mathbf{x}_{t-j}'\Bigr)(\mathbf{b} - \bm{\beta}).
\label{eq:2.10.10}
\end{align}
%
If $\E(\mathbf{x}_t \cdot \varepsilon_{t-j})$, $\E(\mathbf{x}_{t-j} \cdot \varepsilon_t)$, and
$\E(\mathbf{x}_t\,\mathbf{x}_{t-j}')$ are all finite, then the second and the third
terms vanish (converge to zero in probability) because
$\mathbf{b} - \bm{\beta} \pto \mathbf{0}$. Therefore,
\[
\hat{\gamma}_j - \tilde{\gamma}_j \pto 0
\quad (j = 0, 1, 2, \ldots\,),
\]
and thus the difference between $\tilde{\rho}_j$ and $\hat{\rho}_j$, too, vanishes in large samples.
However, $\sqrt{n}$ times the difference does not.
$\sqrt{n}\,\tilde{\rho}_j$ and $\sqrt{n}\,\hat{\rho}_j$ can be written as
%
\begin{equation}
\sqrt{n}\,\tilde{\rho}_j = \frac{\sqrt{n}\,\tilde{\gamma}_j}{\tilde{\gamma}_0}
\quad \text{and} \quad
\sqrt{n}\,\hat{\rho}_j = \frac{\sqrt{n}\,\hat{\gamma}_j}{\hat{\gamma}_0}.
\label{eq:2.10.11}
\end{equation}
%
We have just seen that, for both $\sqrt{n}\,\tilde{\rho}_j$ and $\sqrt{n}\,\hat{\rho}_j$,
the denominator $\pto \sigma^2$. So the difference between $\sqrt{n}\,\tilde{\rho}_j$ and
$\sqrt{n}\,\hat{\rho}_j$ will vanish if the difference between $\sqrt{n}\,\tilde{\gamma}_j$
and $\sqrt{n}\,\hat{\gamma}_j$ does (if you are not convinced, see Review Question~3 below).
Now, by multiplying both sides of \eqref{eq:2.10.10} by $\sqrt{n}$, we obtain
%
\begin{align}
\sqrt{n}\,\hat{\gamma}_j &= \sqrt{n}\,\tilde{\gamma}_j
- \frac{1}{n}\sum_{t=j+1}^{n}
(\mathbf{x}_{t-j}\cdot\varepsilon_t + \mathbf{x}_t\cdot\varepsilon_{t-j})'\,
\sqrt{n}\,(\mathbf{b} - \bm{\beta}) \notag \\
&\quad + \sqrt{n}\,(\mathbf{b} - \bm{\beta})'\Bigl(\frac{1}{n}\sum_{t=j+1}^{n}
\mathbf{x}_t\,\mathbf{x}_{t-j}'\Bigr)(\mathbf{b} - \bm{\beta}).
\label{eq:2.10.12}
\end{align}
%
Because $\sqrt{n}\,(\mathbf{b} - \bm{\beta})$ converges to a random variable (whose distribution is
normal), the third term on the RHS vanishes by Lemma~2.4(b). Regarding the second term,
%
\begin{equation}
\frac{1}{n}\sum_{t=j+1}^{n}
(\mathbf{x}_{t-j}\cdot\varepsilon_t + \mathbf{x}_t\cdot\varepsilon_{t-j})
\pto \E(\mathbf{x}_{t-j}\cdot\varepsilon_t) + \E(\mathbf{x}_t\cdot\varepsilon_{t-j}).
\label{eq:2.10.13}
\end{equation}
%
If the regressors are strictly exogenous in the sense that
$\E(\mathbf{x}_t \cdot \varepsilon_s) = \mathbf{0}$ for all $t$, $s$, then
%
\begin{equation}
\E(\mathbf{x}_{t-j}\cdot\varepsilon_t) + \E(\mathbf{x}_t\cdot\varepsilon_{t-j}) = \mathbf{0},
\label{eq:2.10.14}
\end{equation}
%
and by Lemma~2.4(b) the second term converges to zero in probability. So the
difference between $\sqrt{n}\,\tilde{\rho}_j$ and $\sqrt{n}\,\hat{\rho}_j$ vanishes,
which means that the $Q$ statistic calculated from the regression residuals $\{e_t\}$,
too, is asymptotically chi-squared, and we can use this residual-based $Q$ to test for
serial correlation. If, on the other hand, the regressors are not strictly exogenous,
then there is no guarantee that \eqref{eq:2.10.14} holds. Consequently, the residual-based
$Q$ statistic may not be asymptotically chi-squared.

\subsection*{Testing with Predetermined, but Not Strictly Exogenous, Regressors}

Therefore, when the regressors are not strictly exogenous, we need to modify the
$Q$ statistic to restore its asymptotic distribution. For this purpose, consider two
restrictions:
%
\begin{align}
&\text{(stronger form of predeterminedness)}\quad
\E(\varepsilon_t \mid \varepsilon_{t-1}, \varepsilon_{t-2}, \ldots,
\mathbf{x}_t, \mathbf{x}_{t-1}, \ldots\,) = 0,
\label{eq:2.10.15} \\[6pt]
&\text{(stronger form of conditional homoskedasticity)} \notag \\
&\qquad \E(\varepsilon_t^2 \mid \varepsilon_{t-1}, \varepsilon_{t-2}, \ldots,
\mathbf{x}_t, \mathbf{x}_{t-1}, \ldots\,) = \sigma^2 > 0.
\label{eq:2.10.16}
\end{align}
%
The first condition is just reproducing (2.3.1) (with ``$t$'' now being used as the
subscript). As was shown in Section~2.3, this is stronger than Assumption~2.3 and
implies that $\mathbf{g}_t$ ($\equiv \mathbf{x}_t \cdot \varepsilon_t$) is an m.d.s. Condition~\eqref{eq:2.10.16}, with the conditioning
set including $\mathbf{x}_t$, is obviously stronger than Assumption~2.7 (conditional homoskedasticity). It also is stronger than the own conditional homoskedasticity assumption
in Proposition~\ref{prop:2.9} because the conditioning set includes current and past $\mathbf{x}$ as well
as past $\varepsilon$. The next result shows that under these additional conditions there is an
appropriate modification of the $Q$ statistic.

\begin{proposition}[2.10 (testing for serial correlation with predetermined regressors)]
\label{prop:2.10}
Suppose that Assumptions~2.1, 2.2, 2.4, \eqref{eq:2.10.15}, and \eqref{eq:2.10.16}
are satisfied. Let the sample autocorrelation of the OLS residuals, $\hat{\rho}_j$,
be defined as in \eqref{eq:2.10.8}. Then,
%
\begin{equation}
\sqrt{n}\,\hat{\bm{\gamma}} \dto N(\mathbf{0},\,\sigma^4 \cdot (\mathbf{I}_p - \mathbf{\Phi}))
\quad \text{and} \quad
\sqrt{n}\,\hat{\bm{\rho}} \dto N(\mathbf{0},\,\mathbf{I}_p - \mathbf{\Phi}),
\label{eq:2.10.17}
\end{equation}
%
where
$\hat{\bm{\gamma}} = (\hat{\gamma}_1, \hat{\gamma}_2, \ldots, \hat{\gamma}_p)'$,
$\hat{\bm{\rho}} = (\hat{\rho}_1, \hat{\rho}_2, \ldots, \hat{\rho}_p)'$,
and $\phi_{jk}$ ($\equiv (j,k)$ element of the $p \times p$ matrix $\mathbf{\Phi}$) is given by
%
\begin{equation}
\phi_{jk} = \E(\mathbf{x}_t \cdot \varepsilon_{t-j})'\,
\E(\mathbf{x}_t\mathbf{x}_t')^{-1}\,
\E(\mathbf{x}_t \cdot \varepsilon_{t-k})/\sigma^2.
\label{eq:2.10.18}
\end{equation}
\end{proposition}

The proof, though not very difficult, is relegated to Appendix~2.B. By the Ergodic
Theorem, matrix $\mathbf{\Phi}$ is consistently estimated by its sample counterpart:
%
\begin{equation}
\widehat{\mathbf{\Phi}} \equiv (\hat{\phi}_{jk}), \quad
\hat{\phi}_{jk} \equiv \bar{\bm{\mu}}_j'\,\mathbf{S}_{\mathbf{xx}}^{-1}\,
\bar{\bm{\mu}}_k / s^2
\quad (j, k = 1, 2, \ldots, p),
\label{eq:2.10.19}
\end{equation}
%
where
\[
s^2 \equiv \frac{1}{n - K}\sum_{t=1}^{n} e_t^2, \quad
\bar{\bm{\mu}}_j \equiv \frac{1}{n}\sum_{t=j+1}^{n}\mathbf{x}_t \cdot e_{t-j}.
\]
%
It follows from this and Proposition~\ref{prop:2.10} that
%
\begin{equation}
\text{modified Box--Pierce } Q
\equiv n \cdot \hat{\bm{\rho}}'(\mathbf{I}_p - \widehat{\mathbf{\Phi}})^{-1}\hat{\bm{\rho}}
\dto \chi^2(p).
\label{eq:2.10.20}
\end{equation}
%
As long as the regressors are predetermined and the error is conditionally homoskedastic in the sense of \eqref{eq:2.10.15} and \eqref{eq:2.10.16}, this modified $Q$ statistic can be
used even when the regressors are not strictly exogenous.

\subsection*{An Auxiliary Regression-Based Test}

Although calculating this modified $Q$ statistic is straightforward with matrix-based
software, it is useful to find an asymptotically equivalent statistic that can be calculated from regression packages. For this purpose, consider the following auxiliary
regression:
%
\begin{equation}
\text{regress } e_t \text{ on } \mathbf{x}_t, e_{t-1}, e_{t-2}, \ldots, e_{t-p}.
\label{eq:2.10.21}
\end{equation}
%
To run this auxiliary regression for $t = 1, 2, \ldots, n$, we need data on
$(e_0, e_{-1}, \ldots, e_{-p+1})$. It does not matter asymptotically which particular numbers
to assign to them, but it seems sensible to set them equal to 0, their expected
value.\footnote{Another asymptotically equivalent choice is to run the auxiliary
regression for $t = p+1, p+2, \ldots, n$.}
From this auxiliary regression, we can calculate the $F$ statistic for the hypothesis
that the $p$ coefficients of $e_{t-1}, e_{t-2}, \ldots, e_{t-p}$ are all zero. Given
Proposition~2.5(c), it is only natural to wonder whether $p \cdot F$ is asymptotically
$\chi^2(p)$. This conjecture is indeed true: under the hypothesis of Proposition~\ref{prop:2.10}, the modified $Q$ statistic \eqref{eq:2.10.20} is
asymptotically equivalent to $p \cdot F$ (i.e., the difference between the two converges
to zero in probability as $n \to \infty$), so $p \cdot F$, too, is asymptotically chi-squared
(showing this is an analytical exercise).

This $p \cdot F$ statistic, in turn, is asymptotically equivalent to $nR^2$ from the auxiliary regression. This can be shown as follows. Recall the algebraic result from
Chapter~1 that the $F$-ratio can be calculated from the difference in the sum of
squared residuals between the unrestricted and restricted regressions. The unrestricted regression in the present context is \eqref{eq:2.10.21} while the restricted regression is
%
\begin{equation}
\text{regress } e_t \text{ on } \mathbf{x}_t.
\label{eq:2.10.22}
\end{equation}
%
Therefore, if $SSR_U$ and $SSR_R$ are SSRs from \eqref{eq:2.10.21} and \eqref{eq:2.10.22}, respectively,
we have
%
\begin{equation}
p \cdot F = \frac{SSR_R - SSR_U}{SSR_U/(n - \#\mathbf{x}_t - p)}
= (n - \#\mathbf{x}_t - p)\,\frac{SSR_R - SSR_U}{SSR_U}
\label{eq:2.10.23}
\end{equation}
%
where $\#\mathbf{x}_t$ is the number of variables in $\mathbf{x}_t$. However, since $e_t$ is the residual from
the original regression (a regression of $y_t$ on $\mathbf{x}_t$), the regressors $\mathbf{x}_t$ in \eqref{eq:2.10.22} have
no explanatory power. So $SSR_R = \mathbf{e}'\mathbf{e}$ where $\mathbf{e}$ is the $n$-dimensional vector of
residuals from the original regression, and \eqref{eq:2.10.23} is numerically identical to
\[
(n - \#\mathbf{x}_t - p)\,\frac{R_{uc}^2}{1 - R_{uc}^2},
\]
where $R_{uc}^2$ is the uncentered $R^2$ for the auxiliary regression \eqref{eq:2.10.21}.\footnote{Recall from (1.2.16) that the uncentered $R^2$ for a regression of $\mathbf{y}$ on $\mathbf{x}$ is defined as
$R_{uc}^2 = (\mathbf{y}'\mathbf{y} - \mathbf{e}'\mathbf{e})/(\mathbf{y}'\mathbf{y})$.
In the present context, the $\mathbf{y}'\mathbf{y}$ in this formula is $\mathbf{e}'\mathbf{e}$ while the $\mathbf{e}'\mathbf{e}$ in the formula is $SSR_U$.}
But since the sample mean of $e_t$ is by construction zero (this is because $\mathbf{x}_t$ includes a constant), $R_{uc}^2$ is numerically identical to the $R^2$ for the auxiliary regression. Thus, we
have proved the algebraic result that
\[
p \cdot F = (n - \#\mathbf{x}_t - p)\,\frac{R^2}{1 - R^2},
\]
where $R^2$ is equal to the $R^2$ for the auxiliary regression \eqref{eq:2.10.21}. Solving this
equation for $R^2$ and multiplying both sides by $n$, we obtain
\[
nR^2 = \frac{n}{n - \#\mathbf{x}_t - p} \cdot
\frac{1}{1 + \frac{p \cdot F}{n - \#\mathbf{x}_t - p}} \cdot p \cdot F.
\]
Since $\frac{p \cdot F}{n - \#\mathbf{x}_t - p} \pto 0$ as $n \to \infty$
(an implication of Lemma~2.4(b)), this equation shows that $p \cdot F - nR^2 \pto 0$.
Therefore, $nR^2$ from the auxiliary regression \eqref{eq:2.10.21} is asymptotically
$\chi^2(p)$. The test based on $nR^2$ is called the
\textbf{Breusch--Godfrey test for serial correlation}. When $p = 1$, the test statistic is
asymptotically equivalent to the square of what is known as \textbf{Durbin's $h$}
statistic.\footnote{See Breusch (1978) and Godfrey (1978). The Breusch--Godfrey test
was originally derived as a Lagrange Multiplier test for the case where the error is
normally distributed and $\mathbf{x}_t$ consists of fixed regressors and the lagged dependent
variable. The discussion in the text shows that the test is applicable to the more general
case considered in Proposition~\ref{prop:2.10}.}

\bigskip
\noindent\textsc{Questions for Review}
\begin{enumerate}
\item (Consistency of sample autocovariances)\quad
Show: if $\{z_t\}$ is ergodic stationary,
then the sample autocovariance $\hat{\gamma}_j$ given in \eqref{eq:2.10.1} is consistent for
$\gamma_j$, the population autocovariance.
\textbf{Hint:} $\gamma_j = \E(z_t z_{t-j}) - \E(z_t)\,\E(z_{t-j})$. Rewrite $\hat{\gamma}_j$ as
\[
\hat{\gamma}_j = \frac{1}{n}\sum_{t=j+1}^{n} z_t\,z_{t-j}
- \bar{z}_n\,\frac{1}{n}\sum_{t=j+1}^{n}z_{t-j}
- \bar{z}_n\,\frac{1}{n}\sum_{t=j+1}^{n}z_t
+ \frac{n-j}{n}\,(\bar{z}_n)^2.
\]

\item (Asymptotic equivalence of two $Q$'s)\quad
Prove that the difference between the
Box--Pierce $Q$ and the Ljung--Box $Q$ converges to 0 in probability as $n \to \infty$
(so their asymptotic distributions are the same). \textbf{Hint:} Let
\[
\mathbf{x}_n \equiv \bigl((\sqrt{n}\,\hat{\rho}_1)^2, \ldots,
(\sqrt{n}\,\hat{\rho}_p)^2\bigr)'.
\]
Find a $p$-dimensional vector $\mathbf{a}_n$ such that the difference between two $Q$'s is
$\mathbf{a}_n'\mathbf{x}_n$. Show that $\mathbf{a}_n \to \mathbf{0}$. The asymptotic property to use is Lemma~2.4(b).

\item Consider $\sqrt{n}\,\tilde{\rho}_j$ and $\sqrt{n}\,\hat{\rho}_j$ in \eqref{eq:2.10.11}.
We have shown that $\hat{\gamma}_0 \pto \sigma^2$ and
$\tilde{\gamma}_0 \pto \sigma^2$. By Proposition~\ref{prop:2.9},
$\sqrt{n}\,\tilde{\gamma}_j \dto N(0, \sigma^4)$. Taking these for granted,
show that $\sqrt{n}\,\hat{\rho}_j - \sqrt{n}\,\tilde{\rho}_j \pto 0$
if $\sqrt{n}\,\hat{\gamma}_j - \sqrt{n}\,\tilde{\gamma}_j \pto 0$.
\textbf{Hint:} If $\sqrt{n}\,\hat{\gamma}_j - \sqrt{n}\,\tilde{\gamma}_j \pto 0$,
then by Lemma~2.4(a), $\sqrt{n}\,\hat{\gamma}_j \dto N(0, \sigma^4)$.
\[
\frac{\tilde{\gamma}_j}{\tilde{\gamma}_0}
- \frac{\hat{\gamma}_j}{\hat{\gamma}_0}
= \tilde{\gamma}_j \cdot \Bigl(\frac{1}{\tilde{\gamma}_0}
- \frac{1}{\sigma^2}\Bigr)
- \hat{\gamma}_j \cdot \Bigl(\frac{1}{\hat{\gamma}_0}
- \frac{1}{\sigma^2}\Bigr)
+ \frac{1}{\sigma^2}\cdot(\tilde{\gamma}_j - \hat{\gamma}_j).
\]
Multiply both sides by $\sqrt{n}$ and apply Lemma~2.4(b).
\end{enumerate}
